\section*{第 \ref{ch:g12:mechanical-waves} 章——机械波}

\begin{solution}{exo:g12:mechanical-waves:1}
(a) $v = 320/40 = \qty{8.0}{m/s}$。(b) 没有观众沿
环形看台移动：扰动（起立这一动作）在传播，人群本身并不移动。
(c) $12/1.5 = \qty{8.0}{m/s}$——一致。
\end{solution}

\begin{solution}{exo:g12:mechanical-waves:2}
(a) \omterm{def:g12:mechanical-waves:transverse}{横波}；(b) \omterm{def:g12:mechanical-waves:transverse}{纵波}；(c) \omterm{def:g12:mechanical-waves:transverse}{纵波}；(d) \omterm{def:g12:mechanical-waves:transverse}{横波}；
(e) P 波为\omterm{def:g12:mechanical-waves:transverse}{纵波}，S 波为\omterm{def:g12:mechanical-waves:transverse}{横波}。
\end{solution}

\begin{solution}{exo:g12:mechanical-waves:3}
(a) $d = 340 \times 4.0 \approx \qty{1.4}{km}$.
(b) $1500 \times 4.0 = \qty{6.0}{km}$。(c) 应取所穿过的
介质中的\omterm{def:g12:mechanical-waves:speed}{传播速度}。
\end{solution}

\begin{solution}{exo:g12:mechanical-waves:4}
$\lambda = v/f$：空气中 $340/440 = \qty{0.77}{m}$，
水中 $1500/440 = \qty{3.4}{m}$。\omterm{def:g10:signals-and-waves:frequency}{频率}保持不变（由
波源决定）；$v$ 与 $\lambda$ 均发生变化。
\end{solution}

\begin{solution}{exo:g12:mechanical-waves:5}
$\lambda = \qty{12}{cm}$; $T = \lambda/v = 0.12/0.30 = \qty{0.40}{s}$;
$f = 1/T = \qty{2.5}{Hz}$.
\end{solution}

\begin{solution}{exo:g12:mechanical-waves:6}
(a) $v = 2.4/0.30 = \qty{8.0}{m/s}$。(b) 由第一张照片可知，
波峰还需传播 $\qty{6.0}{m}$，于是 $t = 6.0/8.0 = \qty{0.75}{s}$ 后波峰到达该处。
(c) 绳上的条带先升后降，纯属横向运动，且
始终停留在 $x = \qty{5.0}{m}$ 处：绳子本身并不随波迁移。
\end{solution}

\begin{solution}{exo:g12:mechanical-waves:7}
(a) 接近速率 \qty{4.0}{m/s}，间距 \qty{4.0}{m}：$t = \qty{1.0}{s}$。
(b) 叠加：$3.0 + 2.0 = \qty{5.0}{cm}$。
(c) $3.0 - 2.0 = \qty{1.0}{cm}$。(d) 两个完好的脉冲
穿出，并保持原状继续传播，仿佛它们从未相遇。
\end{solution}

\begin{solution}{exo:g12:mechanical-waves:8}
(a) 细的那根：\omterm{def:g10:inertia:inventory}{张力}相同时，\omterm{def:g10:orders-of-magnitude:unit}{单位}长度的质量越小，
波速越快。(b) 波速增大（\omterm{def:g10:inertia:inventory}{张力}更大）。(c) 否：
波速只取决于介质，与脉冲大小无关。
\end{solution}

\begin{solution}{exo:g12:mechanical-waves:9}
钢轨中：$1800/5900 \approx \qty{0.31}{s}$；空气中：
$1800/340 \approx \qty{5.3}{s}$；时间差 $\approx \qty{5.0}{s}$。
\end{solution}

\begin{solution}{exo:g12:mechanical-waves:10}
取 $\lambda = \qty{24}{cm}$：(a) $2\lambda$，\omterm{def:g12:mechanical-waves:phase}{同相}；
(b) $\lambda/2$，\omterm{def:g12:mechanical-waves:phase}{反相}；(c) $\tfrac32\lambda$，\omterm{def:g12:mechanical-waves:phase}{反相}；
(d) $1.25\lambda$，两者皆非；(e) $\lambda/4$，两者皆非。
\end{solution}

\begin{solution}{exo:g12:mechanical-waves:11}
$\lambda_P = 6.0 \times 2.0 = \qty{12}{km}$;
$\lambda_S = 3.5 \times 2.0 = \qty{7.0}{km}$——约为一栋建筑物尺寸的数百倍。
距离远小于 $\lambda$ 的各点几乎\omterm{def:g12:mechanical-waves:phase}{同相}，
因此整个邻近区域一起运动。
\end{solution}

\begin{solution}{exo:g12:mechanical-waves:12}
(a) $T = \qty{0.50}{s}$, $f = \qty{2.0}{Hz}$.
(b) $\lambda = \qty{0.75}{m}$.
(c) $v = \lambda/T = 0.75/0.50 = \qty{1.5}{m/s}$.
(d) 在 $t = T = \qty{0.50}{s}$ 时。(e) $1.875/0.75 = 2.5$ 个
\omterm{def:g12:mechanical-waves:wavelength}{波长} $= (2 + \tfrac12)\lambda$：\omterm{def:g12:mechanical-waves:phase}{反相}。
\end{solution}

\begin{solution}{exo:g12:mechanical-waves:13}
(a) $T = 80/10 = \qty{8.0}{s}$, $f = \qty{0.125}{Hz}$.
(b) $v = \lambda/T = 120/8.0 = \qty{15}{m/s} = \qty{54}{km/h}$.
(c) $t = 900/15 = \qty{60}{s}$。(d) 否：波不传递物质——
锚定的浮标只是上下起伏。
\end{solution}

\begin{solution}{exo:g12:mechanical-waves:14}
(a) 以接近速率 \qty{6.0}{m/s} 相向运动，间距 \qty{6.0}{m}：$t = \qty{1.0}{s}$；
相反的位移\omterm{def:g10:inertia:compensate}{相互抵消}——绳子此刻呈平直状。
(b) 否：绳上各点仍在运动（横向速度相加，
而非抵消）；\omterm{def:g10:energy-conservation:energy}{能量}全部为\omterm{def:g11:mechanical-energy:kinetic}{动能}。(c) 两个脉冲重新出现，
并保持形状不变地继续传播。(d) 平直的绳子具有速度，因而具有\omterm{def:g10:energy-conservation:energy}{能量}：
波此刻将\omterm{def:g10:energy-conservation:energy}{能量}储存于运动之中，而非形状之中。
\end{solution}

\begin{solution}{exo:g12:mechanical-waves:15}
(a) $2\pi r$：\qty{3.1}{m} 与 \qty{50}{m}——$\times 16$。
(b) \omterm{def:g10:orders-of-magnitude:unit}{单位}长度的\omterm{def:g10:energy-conservation:energy}{能量}降低 $16$ 倍，因此\omterm{def:g10:signals-and-waves:amplitude}{振幅}减小（
\omterm{def:g10:energy-conservation:energy}{能量}随\omterm{def:g10:signals-and-waves:amplitude}{振幅}增大而增大）。(c) 平直\omterm{def:g12:mechanical-waves:wavefront}{波前}不会变长，
只剩吸收作用。(d) 介质是主动的：每位观众
依靠自身肌肉的\omterm{def:g10:energy-conservation:energy}{能量}站立，为波补充\omterm{def:g10:energy-conservation:energy}{能量}。
\end{solution}

\begin{solution}{pb:g12:mechanical-waves:1}
\textbf{1.} P 波沿传播方向压缩：\omterm{def:g12:mechanical-waves:transverse}{纵波}；S 波
侧向振动：\omterm{def:g12:mechanical-waves:transverse}{横波}。

\textbf{2.} 速度较快的 P 波处处先到达。
$t_P = 120/6.0 = \qty{20}{s}$; $t_S = 120/3.5 \approx \qty{34}{s}$.

\textbf{3.} $\Delta t = d/v_S - d/v_P = d\,(1/v_S - 1/v_P)$；对于
\qty{120}{km}: $34.3 - 20 \approx \qty{14}{s}$.

\textbf{4.} 解出 $d$：$d = \dfrac{v_P v_S}{v_P - v_S}\,\Delta t
= \dfrac{6.0 \times 3.5}{2.5}\,\Delta t = \qty{8.4}{km}$，即每秒
时间差对应的距离。

\textbf{5.} $\Delta t \propto d$：某一台站自身的记录即可给出
它到震中的距离，无需知道发震时刻。

\textbf{6.} $d_A = 8.4 \times 25 = \qty{2.1e2}{km}$。

\textbf{7.} 时间差只给出距离，不给出方向：震中位于以 $d_A$ 为半径、
以 A 为圆心的圆周上的某处。

\textbf{8.} $d_B = 8.4 \times 15 = \qty{1.3e2}{km}$；
$d_C = 8.4 \times 32 \approx \qty{2.7e2}{km}$。

\textbf{9.} 两个圆相交于两点，第三个圆择出其一：
三圆的公共交点即为震中。

\textbf{10.} $t_P = 210/6.0 = \qty{35}{s}$：地震发生于
$09{:}14{:}32 - \qty{35}{s} = 09{:}13{:}57$。

\textbf{11.} $\qty{200}{m/s} = \qty{720}{km/h}$——巡航中的
喷气式客机。

\textbf{12.} $t = \num{6.0e5}/200 = \qty{3.0e3}{s} = \qty{50}{min}$。

\textbf{13.} $t_P = 600/6.0 = \qty{100}{s}$。注意：
$3000 - 100 = \qty{2.9e3}{s} \approx \qty{48}{min}$。

\textbf{14.} $\lambda = vT = 200 \times 900 = \qty{1.8e5}{m} =
\qty{180}{km}$：约一米的起伏分布在 \qty{180}{km} 之上——
船只以数分钟的时间爬过一个难以察觉的斜坡。

\textbf{15.} $\lambda = 10 \times 900 = \qty{9.0}{km}$：波
被压缩为原来的二十分之一，\omterm{def:g10:energy-conservation:energy}{能量}集中在更短的长度内，因此其
高度必然增大——海啸在岸边陡然耸起。

\textbf{16.} $2\pi r$：\qty{63}{km}，随后 $\approx \qty{1.3e3}{km}$
——$\times 21$。

\textbf{17.} 同样的\omterm{def:g10:energy-conservation:energy}{能量}分摊在长 $21\times$ 倍的\omterm{def:g12:mechanical-waves:wavefront}{波前}上，
因此即使在无\omterm{def:g11:mechanical-energy:dissipation}{耗散}的地壳中\omterm{def:g10:signals-and-waves:amplitude}{振幅}也会减小；实际地壳
还会附加吸收。

\textbf{18.} S 波是横向剪切波，不能穿过液体：
S 波影区表明（外）核为液态。

\textbf{19.} $\num{1e15}/\num{1.5e9} \approx \num{6.7e5}$：
约七十万列以全速行驶的列车。

\textbf{20.} 地震发生于 $09{:}13{:}57$，位于近海的
三圆交点处，距台站 A \qty{210}{km}。港口的
P 波警报在发震后 \qty{100}{s} 响起，约比
海啸抵达提前 \qty{48}{min}。
\end{solution}
